\documentclass[11pt]{article}

% --- packages (all in a standard TeX Live install) ---
\usepackage[utf8]{inputenc}
\usepackage[T1]{fontenc}
\usepackage{lmodern}
\usepackage[a4paper,margin=1in]{geometry}
\usepackage{amsmath,amssymb}
\usepackage{booktabs}
\usepackage{enumitem}
\usepackage{microtype}
\usepackage{xcolor}
\usepackage[hidelinks,colorlinks=true,linkcolor=teal,urlcolor=teal,citecolor=teal]{hyperref}
\usepackage{abstract}

\newcommand{\valence}{\ensuremath{v}}
\newcommand{\arousal}{\ensuremath{a}}

\title{\textbf{Vibe Shuffle: Browser-Based Affective Music Selection\\
from Facial Valence and Cardiac Arousal}}
\author{Markus Baumann\\
\small Ludwig-Maximilians-Universit\"at M\"unchen\\
\small \texttt{markus.baumann@campus.lmu.de}}
\date{\today}

\begin{document}
\maketitle

\begin{abstract}
\noindent
Vibe Shuffle is a client-only web application that selects music in response to
a listener's affective state, estimated in real time from facial expression and
(optionally) heart-rate variability. All signal processing runs in the browser:
no audio, video, or physiological data leaves the device, and only anonymous
self-report ratings are exported. We describe the affect model---a fusion of a
facial \valence{} channel and a cardiac \arousal{} channel onto Russell's
valence--arousal circumplex~\cite{russell1980}---and the track-ranking rule that
maps the fused state to the next song. We then specify a blinded,
counterbalanced validation protocol that contrasts Vibe Shuffle against a Random
Shuffle baseline, with perceived \emph{mood-fit} as the primary outcome and
\emph{liking} as a control. The system, this paper, and the data schema are
openly available.
\end{abstract}

\section{Introduction}
Listeners routinely choose music to maintain or change how they
feel~\cite{juslin2008}. Most streaming recommenders, however, optimise for
long-run engagement from listening history rather than the listener's
\emph{current} affective state. Vibe Shuffle asks a narrower question: if we
estimate a listener's momentary valence and arousal on-device and pick the next
track to match that state, do listeners perceive the music as fitting their mood
better than a random draw from the same pool?

The design priorities are (i) \textbf{privacy}---raw sensor streams never leave
the browser; (ii) \textbf{falsifiability}---selection is evaluated against a
matched random baseline under blinding; and (iii) \textbf{reproducibility}---the
client, the curated track pool, and the export schema are public.

\section{Affect Model}

\subsection{Signal chain}
The pipeline is fully client-side:
\begin{center}
\small
camera $\rightarrow$ blendshapes $\rightarrow$ valence \quad$\big|$\quad
heart rate $\rightarrow$ HR/RR $\rightarrow$ arousal \quad$\rightarrow$\quad
fusion $\rightarrow$ quadrant $\rightarrow$ track ranking.
\end{center}

\subsection{Valence from facial expression}
A MediaPipe Face Landmarker~\cite{mediapipe2019} produces per-frame blendshape
activations, which are reduced to \emph{happy}, \emph{tense}, and \emph{sad}
scores. A slowly learned personal neutral baseline is subtracted, and the
positive-versus-negative balance is mapped to a continuous valence
$\valence \in [0,1]$. Head and body motion (nose-tip drift plus frame
differencing) feeds a motion channel that nudges the positive score upward and
raises arousal.

\subsection{Arousal from heart rate}
When a Bluetooth heart-rate sensor is connected, HR and RR intervals are parsed
into a $120\,$s personal baseline. Short-window arousal is computed from the
$z$-scored heart rate (higher HR $\rightarrow$ higher arousal) and the
$z$-scored RMSSD (higher short-term variability $\rightarrow$ lower
arousal)~\cite{shaffer2017}:
\begin{equation}
\arousal_{\text{phys}} \;=\; \operatorname{clip}\!\Big(\tfrac{1}{2}
+ \alpha\, z(\mathrm{HR}) - \beta\, z(\mathrm{RMSSD}),\; 0,\; 1\Big).
\end{equation}
SDNN is logged for quality control but excluded from the short-window estimate.

\subsection{Fusion and quadrants}
The fused state is
\begin{equation}
(\valence, \arousal) =
\big(\valence_{\text{face}},\; \arousal_{\text{phys}} + \gamma\, m\big),
\end{equation}
where $m$ is head motion and $\gamma$ its weight. The face sets valence; a usable
ECG sets the arousal base in both directions, and motion adds to arousal on top.
Without a usable ECG, the face/motion channel carries arousal upward only; with
neither face nor ECG, both axes default to $0.5$. Splitting each axis at $0.5$
yields four quadrants---Energetic, Calm, Tense, Melancholic---following the
circumplex~\cite{russell1980,posner2005}.

\subsection{Track selection}
Each track in the fixed pool carries pre-computed valence/arousal coordinates
$(\valence_t, \arousal_t)$. In a Vibe block, the next track minimises distance to
the fused target $(\valence^\star, \arousal^\star)$, preferring the matching
quadrant and penalising recently played tracks:
\begin{equation}
t^\star = \arg\min_{t}\;
\big\lVert (\valence_t, \arousal_t) - (\valence^\star, \arousal^\star) \big\rVert_2
\;+\; \lambda\, r(t),
\end{equation}
where $r(t)$ is a recency penalty. A Random block draws from the same pool with a
per-session seed, so order differs every run while the pool is held constant.

\section{Validation Protocol}
The app implements a blinded, counterbalanced within-subject comparison.
\begin{itemize}[leftmargin=1.4em]
  \item \textbf{Pool}: 100 curated Spotify tracks, 25 per quadrant; identical for
  every participant.
  \item \textbf{Blocks}: two blocks, \texttt{random} and \texttt{vibe}, of 5
  tracks each. Block order is randomised per session and recorded.
  \item \textbf{Listening window}: each track plays for $60\,$s, then the rating
  prompt opens (early rating allowed).
  \item \textbf{Blinding}: the condition is never shown during the session.
\end{itemize}

After each track, two 5-point questions are asked in sequence: \emph{liking}
(``How much do you like this song?'') and \emph{mood-fit} (``How well did it fit
your current mood?''). Mood-fit is the primary outcome; liking is the control,
separating ``did not fit my mood'' from ``I just don't like this song.''

\section{Analysis Plan}
The pre-specified comparison contrasts \texttt{rating\_fit\_1\_to\_5} between
\texttt{vibe} and \texttt{random} blocks, controlling for
\texttt{rating\_like\_1\_to\_5} and accounting for block order. Each completed
session exports one CSV row per track with song features, the detected
$(\valence,\arousal)$, signal-quality flags, and both ratings (see
Table~\ref{tab:schema}).

\begin{table}[t]
\centering
\small
\caption{Selected exported fields (one row per track).}
\label{tab:schema}
\begin{tabular}{@{}ll@{}}
\toprule
Field & Meaning \\
\midrule
\texttt{block\_mode} & \texttt{vibe} or \texttt{random} \\
\texttt{song\_quadrant} & target affect quadrant of the track \\
\texttt{detected\_valence/arousal} & fused listener state \\
\texttt{physiology\_quality} & ECG signal-quality flag \\
\texttt{rating\_fit\_1\_to\_5} & primary outcome (mood-fit) \\
\texttt{rating\_like\_1\_to\_5} & control (liking) \\
\bottomrule
\end{tabular}
\end{table}

\section{Privacy and Limitations}
No backend exists: camera frames, blendshapes, and heart-rate packets are
processed in memory and never transmitted; only the anonymous rating CSV leaves
the browser, at the participant's initiative. The estimator is a lightweight
heuristic, not a clinical affect measure: facial valence is sensitive to
lighting and occlusion, and the arousal channel requires a compatible BLE
sensor. The track pool is small and fixed, so results speak to relative ranking
within that pool rather than to catalogue-scale recommendation.

\section{Availability}
Source code: \url{https://github.com/eybmits/vibe_shuffle}. Live application:
\url{https://eybmits.github.io/vibe_shuffle_site/}. This document is generated
from \texttt{paper/vibe\_shuffle.tex} in the source repository.

\begin{thebibliography}{9}
\bibitem{russell1980} J.~A. Russell. A circumplex model of affect.
\emph{Journal of Personality and Social Psychology}, 39(6):1161--1178, 1980.
\bibitem{posner2005} J.~Posner, J.~A. Russell, and B.~S. Peterson. The
circumplex model of affect. \emph{Development and Psychopathology},
17(3):715--734, 2005.
\bibitem{juslin2008} P.~N. Juslin and D.~V\"astfj\"all. Emotional responses to
music: The need to consider underlying mechanisms. \emph{Behavioral and Brain
Sciences}, 31(5):559--575, 2008.
\bibitem{mediapipe2019} C.~Lugaresi et al. MediaPipe: A framework for building
perception pipelines. \emph{arXiv:1906.08172}, 2019.
\bibitem{shaffer2017} F.~Shaffer and J.~P. Ginsberg. An overview of heart rate
variability metrics and norms. \emph{Frontiers in Public Health}, 5:258, 2017.
\end{thebibliography}

\end{document}
